% ch9_e 第9章 §9.4 (书页 406–418 / p421–p433)
% 块边界判定：
%   起点——书页 406（p421）顶部为承接书页 405 的半句 "find out which $u_m^\mu$ must
%     vanish."（ch9_d 末句停在"……的关键是"），故本块正文首行以 %JOIN% 续接，
%     随后接本页其余正文。
%   终点——书页 418（p433）末段止于完整句 "The gapless spinons appear only at
%     $(\pi/2,\pm\pi/2)$."；书页 419（p434）顶部另起新段 "Knowing the translational
%     symmetry ..."（连同其后的式 (9.6.9)、9.6.2 节与问题 9.6.1–9.6.3），均属下一块。
%     本块自然结束，不设 TAIL 区。
% 蓝本问题记录（均已按图改正）：
%   1) 书页 406 蓝本作 $u_{ij}=\mathrm{i}u_{ij}^0\tau^0+w_{ij}^3\tau^3$，图上为
%      $u_{ij}^3\tau^3$（蓝本 OCR 之误）。
%   2) 式 (9.4.23)、(9.4.24) 蓝本行尾残留 "\tag {9}" 杂项，已删。
%   3) 式 (9.4.28) 蓝本把最后一个规范变换误作 $G_\Gamma$，图上为 $G_T$。
%   4) 9.4.5 节 U1A 标记行蓝本残损（"U1A$_{a_{\eta_x px}}$b$_{\eta_y px}$cd$_y$,"
%      与 "$G_{F_{xy}}$"、"If the corresponding G contains a $\tau^1$, and equal
%      to $\tau^0$ otherwise"），图上为 "U1A$a_{\eta_{xpx}}b_{\eta_{ypx}}cd_{\eta_t}$、
%      $G_{P_{xy}}$"，句子为 "They are equal to $\tau^1$ if the corresponding $G$
%      contains a $\tau^1$, and equal to $\tau^0$ otherwise."，已按图译出。
%   5) 式 (9.6.8) 第二行右列图上印作 $u_{\pmb{i},\pmb{i}+-\pmb{x}+2\pmb{y}}$（多印一
%      个加号），按本意译为 $u_{\pmb{i},\pmb{i}-\pmb{x}+2\pmb{y}}$。
%   6) 书页 418 的 $\Gamma$ 矩阵定义，蓝本排成两栏且重复 $\Gamma_4$；图上为两行
%      （$\Gamma_0,\Gamma_1,\Gamma_2$ / $\Gamma_3,\Gamma_4$），已按图重排。
%   7) 书页 411–412 "contruct"、"we would like to the study"、书页 406–407 跨页
%      "under under"、书页 415 "(9.2.34"（缺右括号）等为原书排印错误，按本意译出。

%JOIN%找出哪些 $u_{\pmb{m}}^{\mu}$ 必须为零。时间反演变换 $G_{T}T$ 与 $180^{\circ}$ 旋转 $G_{P_{x}}P_{x}G_{P_{y}}P_{y}$ 在这里起关键作用。$G_{T}T$ 下的不变性要求 $-u_{\pmb{m}}=\tau^{3}u_{\pmb{m}}\tau^{3}$，因此只有 $u_{\pmb{m}}^{1,2}$ 可以非零。$G_{P_{x}}P_{x}G_{P_{y}}P_{y}$ 下的不变性要求 $u_{\pmb{m}}=u_{-\pmb{m}}$，而由于 $u_{\pmb{m}}^{1,2}$ 为实数且 $u_{-\pmb{m}}=u_{\pmb{m}}^{\dagger}$，$u_{\pmb{m}}=u_{\pmb{m}}^{1}\tau^{1}+u_{\pmb{m}}^{2}\tau^{2}$ 总能满足这一要求，所以 $180^{\circ}$ 旋转不附加任何新条件。其余变换则把不同的 $u_{\pmb{m}}^{1,2}$ 联系起来。$G_{P_{xy}}P_{xy}$ 下的不变性要求 $u_{P_{xy}\pmb{m}}=\tau^{1}u_{\pmb{m}}\tau^{1}$，其中 $P_{xy}\pmb{m}=P_{xy}(m_{x},m_{y})=(m_{y},m_{x})$。于是，最普遍的 Z2A0013 拟设具有如下形式：
\begin{equation*}
\begin{array}{l}
u_{\pmb{m}}=u_{\pmb{m}}^{1}\tau^{1}+u_{\pmb{m}}^{2}\tau^{2}\\[2pt]
u_{P_{x}\pmb{m}}^{1,2}=u_{\pmb{m}}^{1,2},\qquad u_{P_{y}\pmb{m}}^{1,2}=u_{\pmb{m}}^{1,2}\\[2pt]
u_{P_{xy}\pmb{m}}^{1}=u_{\pmb{m}}^{1},\qquad u_{P_{xy}\pmb{m}}^{2}=-u_{\pmb{m}}^{2}
\end{array}
\end{equation*}
把式 (9.2.41) 中的 $(\tau^{3},\tau^{1})$ 换成 $(\tau^{1},\tau^{2})$，并把 $a_{0}$ 认同于 $u_{\pmb{m}=0}$，我们发现 $Z_{2}$ 线性自旋液体的拟设 (9.2.41) 是上述拟设的一个特例。因此，拟设 (9.2.41) 的 PSG 就是 Z2A0013 PSG，相应的自旋液体是一个 Z2A0013 自旋液体。

除了上面研究的 $Z_{2}$ 对称自旋液体之外，还可能存在低能规范群为 $U(1)$ 或 $SU(2)$ 的对称自旋液体。这类 $U(1)$ 对称或 $SU(2)$ 对称的自旋液体由 IGG 为 $U(1)$ 或 $SU(2)$ 的 PSG 分类。我们把这些 PSG 称为 $U(1)$ PSG 与 $SU(2)$ PSG。$U(1)$ PSG 与 $SU(2)$ PSG 已由 Wen (2002c) 算出，9.4.5 节将总结这些结果。

$U(1)$ PSG 中的一个由下式给出：
\begin{equation*}
\begin{array}{rcl}
G_{x}=g_{3}(\theta_{x})\mathrm{i}\tau^{1}, & & G_{y}=g_{3}(\theta_{y})\mathrm{i}\tau^{1},\\[2pt]
G_{P_{x}}=(-)^{i_{x}}g_{3}(\theta_{px}), & & G_{P_{y}}=(-)^{i_{y}}g_{3}(\theta_{py}),\\[2pt]
G_{P_{xy}}=g_{3}(\theta_{pxy})\mathrm{i}\tau^{1}, & & G_{T}=(-)^{\pmb{i}}g_{3}(\theta_{t}).
\end{array}
\tag{9.4.19}
\end{equation*}
其中
\begin{equation*}
g_{a}(\theta)\equiv\mathrm{e}^{\mathrm{i}\theta\tau^{a}}.
\end{equation*}
IGG 由 $\{g_{3}(\phi)\,|\,\phi\in[0,2\pi)\}$ 构成。这样的 PSG 标记为 U1Cn01n。在式 (9.4.27) 中取 $\eta_{xpx}=-1$、$\eta_{ypx}=+1$、$\eta_{pxy}=+1$ 与 $\eta_{t}=-1$ 即得到它。

让我们求出在 U1Cn01n PSG 下不变的最普遍拟设。要在 IGG 下不变，拟设必须对任意 $\theta$ 满足 $g_{3}(\theta)u_{\pmb{ij}}g_{3}^{\dagger}(\theta)=u_{\pmb{ij}}$，因此 $u_{\pmb{ij}}$ 必须具有 $u_{\pmb{ij}}=\mathrm{i}u_{\pmb{ij}}^{0}\tau^{0}+u_{\pmb{ij}}^{3}\tau^{3}$ 的形式。在平移变换 $G_{x}T_{x}$ 与 $G_{y}T_{y}$ 下的不变性要求 $u_{\pmb{ij}}^{0}$ 平移不变，而 $u_{\pmb{ij}}^{3}$ 在两个平移 $T_{x,y}$ 下变号。于是 $u_{\pmb{ij}}$ 具有形式 $u_{\pmb{i},\pmb{i}+\pmb{m}}=\mathrm{i}u_{\pmb{m}}^{0}\tau^{0}+(-)^{i}u_{\pmb{m}}^{3}\tau^{3}$。时间反演变换 $G_{T}T$ 下的不变性要求 $-u_{\pmb{m}}=(-)^{\pmb{m}}u_{\pmb{m}}$。我们发现，若 $\pmb{m}$ 为偶，则 $u_{\pmb{m}}^{0,3}=0$。\footnote{这里 $\pmb{m}$ 为偶（奇），若 $m_{x}+m_{y}$ 为偶（奇）。}$180^{\circ}$ 旋转 $G_{P_{x}}P_{x}G_{P_{y}}P_{y}$ 下的不变性要求 $u_{-\pmb{i},-\pmb{i}-\pmb{m}}=u_{-\pmb{i}-\pmb{m},-\pmb{i}}^{\dagger}=(-)^{\pmb{m}}u_{\pmb{i},\pmb{i}+\pmb{m}}$，这意味着 $-\mathrm{i}u_{\pmb{m}}^{0}\tau^{0}+(-)^{\pmb{i}+\pmb{m}}u_{\pmb{m}}^{3}\tau^{3}=(-)^{\pmb{m}}\big(\mathrm{i}u_{\pmb{m}}^{0}\tau^{0}+(-)^{\pmb{i}}u_{\pmb{m}}^{3}\tau^{3}\big)$。由此得到 $u_{\pmb{m}}^{0}=-(-)^{\pmb{m}}u_{\pmb{m}}^{0}$ 与 $u_{\pmb{m}}^{3}=u_{\pmb{m}}^{3}$，对 $u_{\pmb{m}}^{0,3}$ 没有给出新的约束。厄米关系 $u_{\pmb{i},\pmb{i}+\pmb{m}}^{\dagger}=u_{\pmb{i}+\pmb{m},\pmb{i}}$ 给出 $u_{\pmb{m}}^{0}=-u_{-\pmb{m}}^{0}$ 与 $u_{\pmb{m}}^{3}=(-)^{\pmb{m}}u_{-\pmb{m}}^{3}$。汇总上述条件，并计入 $G_{P_{x}}P_{x}$ 等的作用，我们发现最普遍的 U1Cn01n 拟设具有如下形式：
\begin{equation*}
\begin{array}{l}
u_{\pmb{i},\pmb{i}+\pmb{m}}=\mathrm{i}u_{\pmb{m}}^{0}\tau^{0}+(-)^{i}u_{\pmb{m}}^{3}\tau^{3}\\[2pt]
\qquad u_{\pmb{m}}^{0,3}=0,\ \text{若 }\pmb{m}\text{ 为偶},\qquad u_{\pmb{m}}^{0,3}=-u_{-\pmb{m}}^{0,3}\\[2pt]
\qquad u_{P_{x}\pmb{m}}^{0,3}=(-)^{m_{x}}u_{\pmb{m}}^{0,3},\qquad u_{P_{y}\pmb{m}}^{0,3}=(-)^{m_{y}}u_{\pmb{m}}^{0,3}\\[2pt]
\qquad u_{P_{xy}\pmb{m}}^{0}=u_{\pmb{m}}^{0},\qquad u_{P_{xy}\pmb{m}}^{3}=-u_{\pmb{m}}^{3}
\end{array}
\end{equation*}
$U(1)$ 线性拟设（交错通量相）(9.2.34) 是上述拟设的一个特例。因此，交错通量相 (9.2.34) 是一个 U1Cn01n 态。

我们可以用规范变换 $W_{i}=(\mathrm{i}\tau^{1})^{i}$ 把上述拟设变成平移不变的：
\begin{equation*}
\begin{array}{rcl}
u_{\pmb{i},\pmb{i}+\pmb{m}}=\tilde{u}_{\pmb{m}}^{1}\tau^{1}+\tilde{u}_{\pmb{m}}^{2}\tau^{2} & & \\[2pt]
\tilde{u}_{\pmb{m}}^{1,2}=0,\quad \text{若 }\pmb{m}\text{ 为偶}, & & \\[2pt]
G_{x}(\pmb{i})=g_{3}((-)^{\pmb{i}}\theta_{x}), & & G_{y}(\pmb{i})=g_{3}((-)^{\pmb{i}}\theta_{y}),\\[2pt]
G_{P_{x}}(\pmb{i})=g_{3}((-)^{\pmb{i}}\theta_{px}), & & G_{P_{y}}(\pmb{i})=g_{3}((-)^{\pmb{i}}\theta_{py}),\\[2pt]
G_{P_{xy}}(\pmb{i})=\mathrm{i}g_{3}((-)^{\pmb{i}}\theta_{pxy})\tau^{1}, & & G_{T}(\pmb{i})=(-)^{\pmb{i}}g_{3}((-)^{\pmb{i}}\theta_{t});
\end{array}
\tag{9.4.20}
\end{equation*}
其中我们还列出了变换后 PSG 中的规范变换。IGG 现在具有形式 $IGG=\{g_{3}((-)^{\pmb{i}}\phi)\,|\,\phi\in[0,2\pi)\}$。如果把 $(\tau^{1},\tau^{2})$ 换成 $(\tau^{3},\tau^{1})$，则可以证明上述拟设就 $f$ 费米子而言对应一个 $d$ 波态（见式 (9.2.4) 与 (9.2.30)）。

拟设 (9.4.20) 的自旋子谱由下式给出：
\begin{equation*}
E_{\pm}(\pmb{k})=\pm\sqrt{\Bigg[\sum_{\pmb{m}=\text{奇}}\tilde{u}_{\pmb{m}}^{1}\cos(\pmb{m}\cdot\pmb{k})\Bigg]^{2}+\Bigg[\sum_{\pmb{m}=\text{奇}}\tilde{u}_{\pmb{m}}^{2}\cos(\pmb{m}\cdot\pmb{k})\Bigg]^{2}}
\end{equation*}
我们注意到，对奇的 $\pmb{m}$，$\cos(\pmb{m}\cdot\pmb{k})$ 在 $\pmb{k}=(\pm\frac{\pi}{2},\pm\frac{\pi}{2})$ 处为零，因此 $E_{\pm}(\pmb{k})$ 在 $\pmb{k}=(\pm\frac{\pi}{2},\pm\frac{\pi}{2})$ 处为零。U1Cn01n 自旋液体态在 $\pmb{k}=(\pm\frac{\pi}{2},\pm\frac{\pi}{2})$ 处总有至少四支无能隙自旋子。无能隙自旋子及其晶体动量 $(\pm\frac{\pi}{2},\pm\frac{\pi}{2})$ 是 U1Cn01n 自旋液体态的普适性质。自旋子的无能隙性及其晶体动量受 U1Cn01n PSG 保护。这是一个惊人的结果。作为一个纯玻色模型，U1Cn01n 自旋液体态不破缺任何对称性，然而它却包含一种特殊的量子序，保证了无能隙\emph{费米子}的存在！由于 $\pmb{k}=(\pm\frac{\pi}{2},\pm\frac{\pi}{2})$ 处的无能隙自旋子是 U1Cn01n 态的普适性质，我们可以利用它在实验上探测 U1Cn01n 量子序。

\subsection*{问题 9.4.1}

试通过式 (9.4.3) 找出把 $G_{y}(\pmb{i})=\Theta(\pmb{i})\tau^{0}$ 变为 $G_{y}(\pmb{i})=\tau^{0}$ 的规范变换。这里 $\Theta(\pmb{i})$ 是只取 $\pm1$ 两个值的任意函数。

\subsection*{问题 9.4.2}

\begin{enumerate}
\item 从式 (9.4.14) 求出 Z2Azz13 PSG 的规范变换 $G_{x,y}$、$G_{P_{x},P_{y},P_{xy}}$ 与 $G_{T}$。
\item 求出在 Z2Azz13 PSG 下不变的最普遍拟设。我们将在后面证明，这样的拟设总有无能隙费米子激发。
\end{enumerate}

\subsection*{问题 9.4.3}

证明式 (9.4.20)。（提示：见式 (9.4.3)。）

\subsection{评注：对称 $U(1)$ 与 $SU(2)$ 自旋液体的分类}

\begin{keypoints}
\item 对所有平均场拟设的分类：这些拟设的 PSG 是格点对称性群的 $U(1)$ 或 $SU(2)$ 扩张，即 $PSG/U(1)=SG$ 或 $PSG/SU(2)=SG$。
\end{keypoints}

Wen (2002c) 给出了 $U(1)$ PSG 与 $SU(2)$ PSG 的分类。下面我们只总结结果。刻画平均场对称 $U(1)$ 自旋液体的 PSG 可以分成 U1A、U1B、U1C 与 U1$_{n}^{m}$ 四种类型。有下列 24 个 U1A 型 PSG：
\begin{equation*}
\begin{array}{rcl}
G_{x}=g_{3}(\theta_{x}), & & G_{y}=g_{3}(\theta_{y}),\\[2pt]
G_{P_{x}}=\eta_{ypx}^{i_{y}}g_{3}(\theta_{px}), & & G_{P_{y}}=\eta_{ypx}^{i_{x}}g_{3}(\theta_{py})\\[2pt]
G_{P_{xy}}=g_{3}(\theta_{pxy}),\ g_{3}(\theta_{pxy})\mathrm{i}\tau^{1}, & & G_{T}=\eta_{t}^{\pmb{i}}g_{3}(\theta_{t})\big|_{\eta_{t}=-1},\ \eta_{t}^{\pmb{i}}g_{3}(\theta_{t})\mathrm{i}\tau^{1}
\end{array}
\tag{9.4.21}
\end{equation*}
以及
\begin{equation*}
\begin{array}{rcl}
G_{x}=g_{3}(\theta_{x}), & & G_{y}=g_{3}(\theta_{y}),\\[2pt]
G_{P_{x}}=\eta_{xpx}^{i_{x}}g_{3}(\theta_{px})\mathrm{i}\tau^{1}, & & G_{P_{y}}=\eta_{xpx}^{i_{y}}g_{3}(\theta_{py})\mathrm{i}\tau^{1}\\[2pt]
G_{P_{xy}}=g_{3}(\theta_{pxy}),\ g_{3}(\theta_{pxy})\mathrm{i}\tau^{1}, & & G_{T}=\eta_{t}^{\pmb{i}}g_{3}(\theta_{t})\big|_{\eta_{t}=-1},\ \eta_{t}^{\pmb{i}}g_{3}(\theta_{t})\mathrm{i}\tau^{1}
\end{array}
\tag{9.4.22}
\end{equation*}
我们将用 $\text{U1A}a_{\eta_{xpx}}b_{\eta_{ypx}}cd_{\eta_{t}}$ 来标记这 24 个 PSG。这里 $a$、$b$、$c$ 与 $d$ 分别与 $G_{P_{x}}$、$G_{P_{y}}$、$G_{P_{xy}}$ 及 $G_{T}$ 相联系：若相应的 $G$ 包含一个 $\tau^{1}$，它们就取 $\tau^{1}$，否则取 $\tau^{0}$。一个典型的记号形如 $\text{U1A}\tau_{-}^{1}\tau^{1}\tau^{0}\tau_{-}^{1}$，可缩写为 $\text{U1A}x10x$。

还有下列 24 个 U1B 型 PSG：
\begin{equation*}
\begin{array}{rcl}
G_{x}=(-)^{i_{y}}g_{3}(\theta_{x}), & & G_{y}=g_{3}(\theta_{y}),\\[2pt]
G_{P_{x}}=\eta_{ypx}^{i_{y}}g_{3}(\theta_{px}), & & G_{P_{y}}=\eta_{ypx}^{i_{x}}g_{3}(\theta_{py})\\[2pt]
(-)^{i_{x}i_{y}}G_{P_{xy}}=g_{3}(\theta_{pxy}),\ g_{3}(\theta_{pxy})\mathrm{i}\tau^{1} & & G_{T}=\eta_{t}^{\pmb{i}}g_{3}(\theta_{t})\big|_{\eta_{t}=-1},\ \eta_{t}^{\pmb{i}}g_{3}(\theta_{t})\mathrm{i}\tau^{1}
\end{array}
\tag{9.4.23}
\end{equation*}
以及
\begin{equation*}
\begin{array}{rcl}
G_{x}=(-)^{i_{y}}g_{3}(\theta_{x}), & & G_{y}=g_{3}(\theta_{y}),\\[2pt]
G_{P_{x}}=\eta_{xpx}^{i_{x}}g_{3}(\theta_{px})\mathrm{i}\tau^{1}, & & G_{P_{y}}=\eta_{xpx}^{i_{y}}g_{3}(\theta_{py})\mathrm{i}\tau^{1},\\[2pt]
(-)^{i_{x}i_{y}}G_{P_{xy}}=g_{3}(\theta_{pxy}),\ g_{3}(\theta_{pxy})\mathrm{i}\tau^{1} & & G_{T}=\eta_{t}^{\pmb{i}}g_{3}(\theta_{t})\big|_{\eta_{t}=-1},\ \eta_{t}^{\pmb{i}}g_{3}(\theta_{t})\mathrm{i}\tau^{1}
\end{array}
\tag{9.4.24}
\end{equation*}
我们将用 $\text{U1B}a_{\eta_{xpx}}b_{\eta_{ypx}}cd_{\eta_{t}}$ 来标记这 24 个 PSG。

60 个 U1C 型 PSG 由下式给出：
\begin{equation*}
\begin{array}{rcl}
G_{x}=g_{3}(\theta_{x})\mathrm{i}\tau^{1}, & & G_{y}=g_{3}(\theta_{y})\mathrm{i}\tau^{1},\\[2pt]
G_{P_{x}}=\eta_{xpx}^{i_{x}}\eta_{ypx}^{i_{y}}g_{3}(\theta_{px}), & & G_{P_{y}}=\eta_{ypx}^{i_{x}}\eta_{xpx}^{i_{y}}g_{3}(\theta_{py})\\[2pt]
G_{P_{xy}}=\eta_{pxy}^{i_{x}}g_{3}\big(\eta_{pxy}^{\pmb{i}}\frac{\pi}{4}+\theta_{pxy}\big), & & G_{T}=\eta_{t}^{\pmb{i}}g_{3}(\theta_{t})\big|_{\eta_{t}=-1},\ \eta_{pxy}^{i_{x}}g_{3}(\theta_{t})\mathrm{i}\tau^{1}
\end{array}
\tag{9.4.25}
\end{equation*}
\begin{equation*}
\begin{array}{rcl}
G_{x}=g_{3}(\theta_{x})\mathrm{i}\tau^{1}, & & G_{y}=g_{3}(\theta_{y})\mathrm{i}\tau^{1},\\[2pt]
G_{P_{x}}=\eta_{xpx}^{i_{x}}g_{3}(\theta_{px})\mathrm{i}\tau^{1}, & & G_{P_{y}}=\eta_{xpx}^{i_{y}}\eta_{pxy}^{\pmb{i}}g_{3}(\theta_{py})\mathrm{i}\tau^{1},\\[2pt]
G_{P_{xy}}=\eta_{pxy}^{i_{x}}g_{3}\big(\eta_{pxy}^{\pmb{i}}\frac{\pi}{4}+\theta_{pxy}\big), & & G_{T}=\eta_{t}^{\pmb{i}}g_{3}(\theta_{t})\big|_{\eta_{t}=-1},\ \eta_{pxy}^{i_{x}}\eta_{t}^{\pmb{i}}g_{3}(\theta_{t})\mathrm{i}\tau^{1}
\end{array}
\tag{9.4.26}
\end{equation*}
\begin{equation*}
\begin{array}{rcl}
G_{x}=g_{3}(\theta_{x})\mathrm{i}\tau^{1}, & & G_{y}=g_{3}(\theta_{y})\mathrm{i}\tau^{1},\\[2pt]
G_{P_{x}}=\eta_{xpx}^{i_{x}}\eta_{ypx}^{i_{y}}g_{3}(\theta_{px}), & & G_{P_{y}}=\eta_{ypx}^{i_{x}}\eta_{xpx}^{i_{y}}g_{3}(\theta_{py}),\\[2pt]
G_{P_{xy}}=g_{3}(\theta_{pxy})\mathrm{i}\tau^{1}, & & G_{T}=\eta_{t}^{\pmb{i}}g_{3}(\theta_{t})\big|_{\eta_{t}=-1}
\end{array}
\tag{9.4.27}
\end{equation*}
\begin{equation*}
\begin{array}{rcl}
G_{x}=g_{3}(\theta_{x})\mathrm{i}\tau^{1}, & & G_{y}=g_{3}(\theta_{y})\mathrm{i}\tau^{1},\\[2pt]
G_{P_{x}}=\eta_{xpx}^{i_{x}}\eta_{ypx}^{i_{y}}g_{3}(\theta_{px}), & & G_{P_{y}}=\eta_{ypx}^{i_{x}}\eta_{xpx}^{i_{y}}g_{3}(\theta_{py}),\\[2pt]
G_{P_{xy}}=g_{3}\big(\eta_{pxy}^{\pmb{i}}\frac{\pi}{4}+\theta_{pxy}\big)\mathrm{i}\tau^{1}, & & G_{T}=\eta_{pxy}^{i_{x}}\eta_{t}^{\pmb{i}}g_{3}(\theta_{t})\mathrm{i}\tau^{1}
\end{array}
\tag{9.4.28}
\end{equation*}
以及
\begin{equation*}
\begin{array}{rcl}
G_{x}=g_{3}(\theta_{x})\mathrm{i}\tau^{1}, & & G_{y}=g_{3}(\theta_{y})\mathrm{i}\tau^{1},\\[2pt]
G_{P_{x}}=\eta_{xpx}^{i_{x}}g_{3}(\theta_{px})\mathrm{i}\tau^{1}, & & G_{P_{y}}=\eta_{xpx}^{i_{y}}\eta_{pxy}^{\pmb{i}}g_{3}(\theta_{py})\mathrm{i}\tau^{1},\\[2pt]
G_{P_{xy}}=g_{3}\big(\eta_{pxy}^{\pmb{i}}\frac{\pi}{4}+\theta_{pxy}\big)\mathrm{i}\tau^{1}, & & G_{T}=\eta_{t}^{\pmb{i}}g_{3}(\theta_{t})\big|_{\eta_{t}=-1},\ \eta_{t}^{\pmb{i}}\eta_{pxy}^{i_{x}}g_{3}(\theta_{t})\mathrm{i}\tau^{1}.
\end{array}
\tag{9.4.29}
\end{equation*}
它们将用 $\text{U1C}a_{\eta_{xpx}}b_{\eta_{ypx}}c_{\eta_{pxy}}d_{\eta_{t}}$ 来标记。

U1$_{n}^{m}$ 型 PSG 尚未被分类。不过我们知道，对每一个有理数 $m/n\in(0,1)$，至少存在一个 U1$_{n}^{m}$ 型的平均场对称自旋液体，其拟设由下式给出：
\begin{equation*}
u_{\pmb{i},\pmb{i}+\pmb{x}}=\chi\tau^{3},\qquad u_{\pmb{i},\pmb{i}+\pmb{y}}=\chi g_{3}\big(\frac{m\pi}{n}i_{x}\big)\tau^{3}
\tag{9.4.30}
\end{equation*}
每个方格穿过 $\pi m/n$ 的通量。因此存在无穷多种 U1$_{n}^{m}$ 自旋液体。

我们想指出，上面 108 个 U1A、U1B 与 U1C PSG 都是代数 PSG，它们只是所有可能的代数 $U(1)$ PSG 的一个子集。不过，它们确实包含了 U1A、U1B 与 U1C 型的所有不变 $U(1)$ PSG。我们发现 108 个 PSG 中有 46 个同时也是不变 PSG，因此共有 46 种 U1A、U1B、U1C 型的平均场 $U(1)$ 自旋液体。它们的拟设与标记见 Wen (2002c)。

为了给对称 $SU(2)$ 自旋液体分类，我们发现共有 8 个不同的 $SU(2)$ PSG，由下式给出：
\begin{equation*}
\begin{array}{rcl}
G_{x}(\pmb{i})=g_{x}, & & G_{y}(\pmb{i})=g_{y},\\[2pt]
G_{P_{x}}(\pmb{i})=\eta_{xpx}^{i_{x}}\eta_{xpy}^{i_{y}}g_{P_{x}}, & & G_{P_{y}}(\pmb{i})=\eta_{xpy}^{i_{x}}\eta_{xpx}^{i_{y}}g_{P_{y}}\\[2pt]
G_{P_{xy}}(\pmb{i})=g_{P_{xy}}, & & G_{T}(\pmb{i})=(-)^{\pmb{i}}g_{T},
\end{array}
\tag{9.4.31}
\end{equation*}
以及
\begin{equation*}
\begin{array}{rcl}
G_{x}(\pmb{i})=(-)^{i_{y}}g_{x}, & & G_{y}(\pmb{i})=g_{y},\\[2pt]
G_{P_{x}}(\pmb{i})=\eta_{xpx}^{i_{x}}\eta_{xpy}^{i_{y}}g_{P_{x}}, & & G_{P_{y}}(\pmb{i})=\eta_{xpy}^{i_{x}}\eta_{xpx}^{i_{y}}g_{P_{y}}\\[2pt]
G_{P_{xy}}(\pmb{i})=(-)^{i_{x}i_{y}}g_{P_{xy}}, & & G_{T}(\pmb{i})=(-)^{\pmb{i}}g_{T}
\end{array}
\tag{9.4.32}
\end{equation*}
其中各个 $g$ 都是 $SU(2)$ 中的矩阵。我们引入如下记号：
\begin{equation*}
\begin{array}{l}
\text{SU2A}\tau_{\eta_{xpx}}^{0}\tau_{\eta_{xpy}}^{0}\\[2pt]
\text{SU2B}\tau_{\eta_{xpx}}^{0}\tau_{\eta_{xpy}}^{0}
\end{array}
\tag{9.4.33}
\end{equation*}
用来表示上述 8 个 PSG：$\text{SU2A}\tau_{\eta_{xpx}}^{0}\tau_{\eta_{xpy}}^{0}$ 对应式 (9.4.31)，$\text{SU2B}\tau_{\eta_{xpx}}^{0}\tau_{\eta_{xpy}}^{0}$ 对应式 (9.4.32)。我们发现，8 个 $SU(2)$ PSG 中只有 4 个——即 $\text{SU2A}[n0,0n]$ 与 $\text{SU2B}[n0,0n]$——能给出 $SU(2)$ 对称自旋液体。SU2A$n0$ 态就是均匀 RVB 态 (9.2.32)，SU2B$n0$ 态就是 $\pi$ 通量态 (9.2.33)。另外两个 $SU(2)$ 自旋液体一个以 SU2A$0n$ 标记，即
\begin{equation*}
\begin{array}{rcl}
u_{\pmb{i},\pmb{i}+2\pmb{x}+\pmb{y}}=+\mathrm{i}\chi\tau^{0}, & & u_{\pmb{i},\pmb{i}-2\pmb{x}+\pmb{y}}=-\mathrm{i}\chi\tau^{0},\\[2pt]
u_{\pmb{i},\pmb{i}+\pmb{x}+2\pmb{y}}=+\mathrm{i}\chi\tau^{0}, & & u_{\pmb{i},\pmb{i}-\pmb{x}+2\pmb{y}}=+\mathrm{i}\chi\tau^{0},
\end{array}
\tag{9.4.34}
\end{equation*}
另一个以 SU2B$0n$ 标记，即
\begin{equation*}
\begin{array}{rcl}
u_{\pmb{i},\pmb{i}+2\pmb{x}+\pmb{y}}=+\mathrm{i}(-)^{i_{x}}\chi\tau^{0}, & & u_{\pmb{i},\pmb{i}-2\pmb{x}+\pmb{y}}=-\mathrm{i}(-)^{i_{x}}\chi\tau^{0},\\[2pt]
u_{\pmb{i},\pmb{i}+\pmb{x}+2\pmb{y}}=+\mathrm{i}\chi\tau^{0}, & & u_{\pmb{i},\pmb{i}-\pmb{x}+2\pmb{y}}=+\mathrm{i}\chi\tau^{0}
\end{array}
\tag{9.4.35}
\end{equation*}

上述结果给出了平均场层面对 $U(1)$ 与 $SU(2)$ 对称自旋液体的分类。我们想指出，$U(1)$ 与 $SU(2)$ 平均场态并不是稳定的平均场态，其中一些是边缘的。因此，与 $Z_{2}$ 平均场态相比，$U(1)$ 与 $SU(2)$ 平均场态是否对应于自旋-1/2 模型某种物理的 $U(1)$ 或 $SU(2)$ 对称自旋液体，就更不清楚了。物理的 $U(1)$ 或 $SU(2)$ 态可能只存在于大 $N$ 极限中（见 9.8 节）。

\section{无对称性破缺的连续相变}

\begin{keypoints}
\item 量子相之间的连续相变由下述原理支配：令 $PSG_{1}$ 与 $PSG_{2}$ 为相变两侧两个量子相的 PSG，$PSG_{cr}$ 为描述量子临界态（quantum critical state）的 PSG，则 $PSG_{1}\subseteq PSG_{cr}$ 且 $PSG_{2}\subseteq PSG_{cr}$。
\item 上述原理既适用于对称性破缺相变，也适用于不改变任何对称性的相变。
\end{keypoints}

在对平均场对称自旋液体进行了分类之后，我们想知道这些对称自旋液体彼此之间有什么关系。特别地，我们想知道哪些自旋液体可以通过\emph{连续}相变相互转变。在平均场层面上，这个问题可以完全用 PSG 来解决。其想法是：PSG 正是平均场态的“对称性”群。平均场相变可以用 PSG 的改变来描述。正如对称性破缺相变一样，PSG 为 $PSG_{1}$ 的平均场态能够通过一个连续的（平均场）相变变为 PSG 为 $PSG_{2}$ 的平均场态，当且仅当 $PSG_{2}\subset PSG_{1}$ 或 $PSG_{1}\subset PSG_{2}$。

为了理解这一结果，让我们假定由 $PSG_{1}$ 描述的平均场态有一个拟设 $u_{\pmb{ij}}$，它的邻居有一个拟设 $u_{\pmb{ij}}+\delta u_{\pmb{ij}}$，这里 $\delta u_{\pmb{ij}}$ 是一个小微扰。设该微扰把 PSG 改变为另一个不同的 PSG——$PSG_{2}$。由于 $\delta u_{\pmb{ij}}$ 是强度未定的微扰，要使 $u_{\pmb{ij}}+\delta u_{\pmb{ij}}$ 在 $PSG_{2}$ 下不变，$u_{\pmb{ij}}$ 与 $\delta u_{\pmb{ij}}$ 就都必须在 $PSG_{2}$ 下不变。因此 $PSG_{2}\subset PSG_{1}$。

利用上述结果，我们可以按下列步骤求出某些重要对称自旋液体邻域中的对称自旋液体。我们从一个 PSG 为 $PSG_{1}$ 的对称自旋液体 $u_{\pmb{ij}}$ 出发，这里 $PSG_{1}$ 是对称性群 $SG$ 被 $IGG_{1}$ 的扩张，即 $PSG_{1}/IGG_{1}=SG$。然后我们找出 $PSG_{1}$ 的所有作为同一对称性群 $SG$ 的扩张的子群，把这些子群记为 $PSG_{2}$。那么 $PSG_{2}$ 必须有一个正规子群 $IGG_{2}$ 使得 $PSG_{2}/IGG_{2}=SG$。利用这些子群，我们可以构造出所有具有相同对称性的近邻平均场态。

经过 Wen (2002c) 中的冗长计算，人们求出了式 (9.2.34) 的 $U(1)$ 线性态 U1Cn01n、式 (9.2.32) 的 $SU(2)$ 无能隙态 SU2An0 以及式 (9.2.33) 的 $SU(2)$ 线性态 SU2Bn0 周围的所有平均场对称自旋液体。例如，已经证明，在平均场层面上，$U(1)$ 线性自旋液体 U1Cn01n 可以连续变为 8 个不同的 $Z_{2}$ 自旋液体：Z2A0013、Z2Azz13、Z2A001n、Z2Azz1n、Z2B0013、Z2Bzz13、Z2B001n 与 Z2Bzz1n。

让我们讨论一个简单的例子来演示上述结果。拟设
\begin{equation*}
\begin{array}{rcl}
u_{\pmb{i},\pmb{i}+\pmb{x}}=\chi\tau^{1}-\eta\tau^{2}, & & u_{\pmb{i},\pmb{i}+\pmb{y}}=\chi\tau^{1}+\eta\tau^{2},\\[2pt]
u_{\pmb{i},\pmb{i}+\pmb{x}+\pmb{y}}=-\gamma\tau^{1}, & & u_{\pmb{i},\pmb{i}-\pmb{x}+\pmb{y}}=+\gamma\tau^{1},\qquad a_{0}^{l}=0.
\end{array}
\tag{9.5.1}
\end{equation*}
在 $\chi\neq\eta$ 时，$\gamma=0$ 的情形描述 U1Cn01n 自旋液体（交错通量态）。也就是说，该拟设在式 (9.4.20) 所给规范变换之后跟随的对称性变换下不变。当 $\gamma\neq0$ 时，拟设生成非共线的 $SU(2)$ 通量，把 $U(1)$ 规范结构破缺到 $Z_{2}$ 规范结构，拟设此时描述一个 $Z_{2}$ 态。对非零的 $\gamma$，拟设不再在式 (9.4.20) 所列 U1Cn01n 规范变换之后跟随的对称性变换下不变，而只在 U1Cn01n 规范变换的一个子集之后跟随的对称性变换下不变：
\begin{equation*}
\begin{array}{rcl}
G_{x}(\pmb{i})=g_{3}(0)=\tau^{0}, & & G_{y}(\pmb{i})=g_{3}(0)=\tau^{0},\\[2pt]
G_{P_{x}}(\pmb{i})=g_{3}((-)^{\pmb{i}}\pi/2)=\mathrm{i}(-)^{\pmb{i}}\tau^{3}, & & G_{P_{y}}(\pmb{i})=g_{3}((-)^{\pmb{i}}\pi/2)=\mathrm{i}(-)^{\pmb{i}}\tau^{3},\\[2pt]
G_{P_{xy}}(\pmb{i})=\mathrm{i}g_{3}(0)\tau^{1}=\mathrm{i}\tau^{1}, & & G_{T}(\pmb{i})=(-)^{\pmb{i}}g_{3}((-)^{\pmb{i}}\pi/2)=\mathrm{i}\tau^{3}.
\end{array}
\end{equation*}
上述规范变换之后跟随的对称性变换给出 Z2Azz13 PSG。于是，当 $\gamma$ 从零变为非零时，U1Cn01n 态连续地变为 Z2Azz13 态。

我们想强调，上述关于连续相变的结果只在平均场层面上有效。有些平均场结果能在量子涨落下保存下来，有些则不能。要知道哪些平均场结果在平均场理论之外仍然有效，需要逐例研究（Mudry and Fradkin, 1994）。

量子涨落的一个可能效果是使某些平均场态失稳。我们可以假定这种失稳来自量子涨落产生的某些相关微扰。在这一假定下，一旦计入量子涨落，一些平均场的稳定不动点对真实物理系统而言就变成了不稳定不动点。于是，平均场相图中的一些相可能收缩成临界线（critical line），这些线代表两个稳定量子相之间相变点上的临界态（critical state）（见图 9.7）。这一图像引出了本节开头要点中所述的关于支配连续量子相变的原理的猜想。

%FIG{9.7}: {(a) 由 $PSG_{1,2,3}$ 刻画三个相的平均场相图，它们之间的平均场相变是连续的。PSG 满足 $PSG_{1}\subset PSG_{3}$ 与 $PSG_{2}\subset PSG_{3}$。$PSG_{1,2,3}$ 的一种可能取法是 $PSG_{1}=$Z2A0013、$PSG_{2}=$Z2Azz13、$PSG_{3}=$U1Cn01n。(b) 计入量子涨落后相应的物理相图。这里我们假定，计入量子涨落之后，态 $PSG_{3}$ 变成一个只带一个相关微扰的不稳定不动点。平均场态 $PSG_{3}$ 收缩成一条临界线，它描述两个态 $PSG_{1}$ 与 $PSG_{2}$ 之间的相变。交错通量态 U1Cn01n 如果不稳定，可能作为一个临界态出现。}

我们想强调，上述所有自旋液体都具有相同的对称性。因此，它们之间的连续相变（如果存在的话）代表了一类新的量子连续相变——不改变任何对称性的相变（Wegner, 1971; Kosterlitz and Thouless, 1973; Dasgupta and Halperin, 1981; Wen and Wu, 1993; Chen \emph{et al.}, 1993; Senthil \emph{et al.}, 1999; Read and Green, 2000; Wen, 2000）。

\subsection*{问题 9.5.1}

拟设
\begin{equation*}
\begin{array}{rcl}
u_{\pmb{i},\pmb{i}+\pmb{x}}=\chi\tau^{1}-\eta\tau^{2}, & & u_{\pmb{i},\pmb{i}+\pmb{y}}=\chi\tau^{1}+\eta\tau^{2},\qquad a_{0}^{1,2,3}=0,\\[2pt]
u_{\pmb{i},\pmb{i}+2\pmb{x}+\pmb{y}}=u_{\pmb{i},\pmb{i}-\pmb{x}+2\pmb{y}}=+\lambda\tau^{3}, & & u_{\pmb{i},\pmb{i}+2\pmb{x}-\pmb{y}}=u_{\pmb{i},\pmb{i}+\pmb{x}+2\pmb{y}}=-\lambda\tau^{3},
\end{array}
\end{equation*}
在 $\chi\neq\eta$ 时，$\lambda=0$ 的情形描述 U1Cn01n 自旋液体（交错通量态）。
\begin{enumerate}
\item 证明：当 $\lambda\neq0$ 时，该拟设描述一个 $Z_{2}$ 自旋液体。
\item 找出保持该拟设不变的 U1Cn01n PSG 的子群。
\item 找出上述 $Z_{2}$ 拟设的标记。
\end{enumerate}

\section{对称自旋液体动物园}

\begin{keypoints}
\item 交错通量态邻域中 8 个 $Z_{2}$ 自旋液体的物理性质。
\end{keypoints}

本节中，我们想研究 9.4 节所分类的对称自旋液体的物理性质。我们一直用投影对称性来刻画不同的平均场态。然而，从一个平均场态的 PSG 直接看出它的物理性质并不容易。为了找出平均场态的物理性质，我们需要构造在相应 PSG 下不变的显式拟设。

然而，把几百个对称自旋液体一一列出绝非易事，更不用说逐一构造它们的拟设并研究其理论性质了。我们这里想做的是研究一些重要的自旋液体；可是，如何判定一个自旋液体的重要性呢？

在对高 $T_{c}$ 超导体的研究中，人们发现 $SU(2)$ 线性态 SU2Bn0（$\pi$ 通量态）、$U(1)$ 线性态 U1Cn01n（交错通量/$d$ 波态）与 $SU(2)$ 无能隙态 SU2An0（均匀 RVB 态）是重要的。它们与高 $T_{c}$ 超导体中观察到的某些相密切相关。$SU(2)$ 线性态、$U(1)$ 线性态与 $SU(2)$ 无能隙态分别再现了未掺杂、欠掺杂与过掺杂样品中观察到的电子谱函数。然而在理论上，由于 $U(1)$ 或 $SU(2)$ 规范涨落，这些自旋液体在低能下是不稳定的。这些态可能会变为它们邻域中更稳定的自旋液体。这促使我们研究 $SU(2)$ 线性态、$U(1)$ 线性态与 $SU(2)$ 无能隙态邻域中这些更稳定的自旋液体。为了进一步缩小范围，本节将主要研究 $U(1)$ 线性态 U1Cn01n 邻域中的自旋液体。\footnote{我们想指出，这里我们只研究对称自旋液体。U1Cn01n 自旋液体也可能变为其他一些破缺某些对称性的态。这样的对称性破缺相变实际上已经在高 $T_{c}$ 超导体中被观察到（如向反铁磁态、$d$ 波超导态与条纹态的转变）。}

\subsection{$U(1)$ 线性自旋液体周围的对称自旋液体}

U1Cn01n 自旋液体 (9.2.34) 可以连续变为 8 个不同的自旋液体，它们把 $U(1)$ 规范结构破缺到 $Z_{2}$ 规范结构。这 8 个自旋液体尽管都具有相同的对称性，却由不同的 PSG 标记。下面我们将更详细地研究这 8 个 $Z_{2}$ 自旋液体。特别地，我们想求出它们之中的自旋子谱。

第一个以 Z2A0013 标记，具有如下形式：
\begin{equation*}
\begin{array}{rcl}
u_{\pmb{i},\pmb{i}+\pmb{x}}=\chi\tau^{1}-\eta\tau^{2}, & & u_{\pmb{i},\pmb{i}+\pmb{y}}=\chi\tau^{1}+\eta\tau^{2},\\[2pt]
u_{\pmb{i},\pmb{i}+\pmb{x}+\pmb{y}}=+\gamma\tau^{1}, & & u_{\pmb{i},\pmb{i}-\pmb{x}+\pmb{y}}=+\gamma\tau^{1},\\[2pt]
a_{0}^{1}\neq0,\qquad a_{0}^{2,3}=0. & &
\end{array}
\tag{9.6.1}
\end{equation*}
它与拟设 (9.2.41) 具有相同的量子序。标记 Z2A0013 告诉我们该拟设的 PSG——即它的“对称性”群。提醒一下，对称自旋液体的 PSG 由对称性变换与规范变换的复合 $\{G_{0},\allowbreak\,G_{x}T_{x},\allowbreak\,G_{y}T_{y},\allowbreak\,G_{P_{x}}P_{x},\allowbreak\,G_{P_{y}}P_{y},\allowbreak\,G_{P_{xy}}P_{xy},\allowbreak\,G_{T}T\}$ 生成。规范变换 $G_{0}$ 生成该 PSG 的 IGG。标记中的“Z2”告诉我们 IGG 为 $Z_{2}$ 且 $G_{0}(\pmb{i})=(-)$；“A”告诉我们 $G_{x}(\pmb{i})=G_{y}(\pmb{i})=1$；接下来的“00”意味着 $G_{P_{x}}(\pmb{i})=1$、$G_{P_{y}}(\pmb{i})=1$；“1”意味着 $G_{P_{xy}}(\pmb{i})=\mathrm{i}\tau^{1}$，而“3”意味着 $G_{T}(\pmb{i})=\mathrm{i}\tau^{3}$。第二个拟设以 Z2Azz13 标记，具有如下形式：
\begin{equation*}
\begin{array}{rcl}
u_{\pmb{i},\pmb{i}+\pmb{x}}=\chi\tau^{1}-\eta\tau^{2}, & & u_{\pmb{i},\pmb{i}+\pmb{y}}=\chi\tau^{1}+\eta\tau^{2},\\[2pt]
u_{\pmb{i},\pmb{i}+\pmb{x}+\pmb{y}}=-\gamma\tau^{1}, & & u_{\pmb{i},\pmb{i}-\pmb{x}+\pmb{y}}=+\gamma\tau^{1},\\[2pt]
u_{\pmb{i},\pmb{i}+2\pmb{x}}=u_{\pmb{i},\pmb{i}+2\pmb{y}}=0,\qquad a_{0}^{1,2,3}=0. & &
\end{array}
\tag{9.6.2}
\end{equation*}
现在标记中的“zz”告诉我们 $\eta_{xpx}=\eta_{xpy}=-1$，$g_{P_{x}}=g_{P_{y}}=\mathrm{i}\tau^{3}$，于是 $G_{P_{x}}(\pmb{i})=(-)^{i_{x}+i_{y}}\tau^{3}$、$G_{P_{y}}(\pmb{i})=(-)^{i_{x}+i_{y}}\tau^{3}$（见式 (9.4.14)）。第三个拟设以 Z2A001n 标记，具有如下形式：
\begin{equation*}
\begin{array}{rcl}
a_{0}^{l}=0, & & \\[2pt]
u_{\pmb{i},\pmb{i}+\pmb{x}}=\chi\tau^{1}+\eta\tau^{2}, & & u_{\pmb{i},\pmb{i}+\pmb{y}}=\chi\tau^{1}-\eta\tau^{2}\\[2pt]
u_{\pmb{i},\pmb{i}+2\pmb{x}+\pmb{y}}=\chi_{1}\tau^{1}+\eta_{1}\tau^{2}+\lambda_{1}\tau^{3}, & & u_{\pmb{i},\pmb{i}-\pmb{x}+2\pmb{y}}=\chi_{1}\tau^{1}-\eta_{1}\tau^{2}-\lambda_{1}\tau^{3},\\[2pt]
u_{\pmb{i},\pmb{i}+2\pmb{x}-\pmb{y}}=\chi_{1}\tau^{1}+\eta_{1}\tau^{2}+\lambda_{1}\tau^{3}, & & u_{\pmb{i},\pmb{i}+\pmb{x}+2\pmb{y}}=\chi_{1}\tau^{1}-\eta_{1}\tau^{2}-\lambda_{1}\tau^{3}.
\end{array}
\tag{9.6.3}
\end{equation*}
如果你足够细心，就会发现标记 Z2A001n 并不出现在 9.4.3 节所分类的 196 个 $Z_{2}$ 自旋液体的名单之中（见表 9.1）。不过，以 Z2A001n 标记的 PSG 与以 Z2A003n 标记的 PSG 是规范等价的，而 Z2A003n 确实出现在名单中（见表 9.1 第二行）。下面我们将把上述自旋液体称为 Z2A003n 态。第四个拟设以 Z2Azz1n 标记，具有如下形式：
\begin{equation*}
\begin{array}{rcl}
a_{0}^{l}=0, & & \\[2pt]
u_{\pmb{i},\pmb{i}+\pmb{x}}=\chi\tau^{1}+\eta\tau^{2}, & & u_{\pmb{i},\pmb{i}+\pmb{y}}=\chi\tau^{1}-\eta\tau^{2}\\[2pt]
u_{\pmb{i},\pmb{i}+2\pmb{x}+\pmb{y}}=\chi_{1}\tau^{1}+\eta_{1}\tau^{2}+\lambda\tau^{3}, & & u_{\pmb{i},\pmb{i}-\pmb{x}+2\pmb{y}}=\chi_{1}\tau^{1}-\eta_{1}\tau^{2}+\lambda\tau^{3},\\[2pt]
u_{\pmb{i},\pmb{i}+2\pmb{x}-\pmb{y}}=\chi_{1}\tau^{1}+\eta_{1}\tau^{2}-\lambda\tau^{3}, & & u_{\pmb{i},\pmb{i}+\pmb{x}+2\pmb{y}}=\chi_{1}\tau^{1}-\eta_{1}\tau^{2}-\lambda\tau^{3}.
\end{array}
\tag{9.6.4}
\end{equation*}

上面四个拟设具有平移不变性。接下来四个 $Z_{2}$ 拟设不具有平移不变性，因为它们都是 Z2B 型的。（不过，投影之后它们仍然描述平移对称的自旋液体。）这些 $Z_{2}$ 自旋液体如下：

Z2B0013：
\begin{equation*}
\begin{array}{rcl}
u_{\pmb{i},\pmb{i}+\pmb{x}}=\chi\tau^{1}-\eta\tau^{2}, & & u_{\pmb{i},\pmb{i}+\pmb{y}}=(-)^{i_{x}}(\chi\tau^{1}+\eta\tau^{2}),\\[2pt]
u_{\pmb{i},\pmb{i}+2\pmb{x}}=-\gamma_{2}\tau^{1}+\lambda_{2}\tau^{2}, & & u_{\pmb{i},\pmb{i}+2\pmb{y}}=-\gamma_{2}\tau^{1}-\lambda_{2}\tau^{2},\\[2pt]
a_{0}^{1}\neq0,\qquad a_{0}^{2,3}=0, & &
\end{array}
\tag{9.6.5}
\end{equation*}

Z2Bzz13：
\begin{equation*}
\begin{array}{rcl}
u_{\pmb{i},\pmb{i}+\pmb{x}}=\chi\tau^{1}-\eta\tau^{2}, & & u_{\pmb{i},\pmb{i}+\pmb{y}}=(-)^{i_{x}}(\chi\tau^{1}+\eta\tau^{2}),\\[2pt]
u_{\pmb{i},\pmb{i}+2\pmb{x}+2\pmb{y}}=-\gamma_{1}\tau^{1}, & & u_{\pmb{i},\pmb{i}-2\pmb{x}+2\pmb{y}}=\gamma_{1}\tau^{1},\\[2pt]
a_{0}^{1,2,3}=0, & &
\end{array}
\tag{9.6.6}
\end{equation*}

Z2B001n：
\begin{equation*}
\begin{array}{rcl}
u_{\pmb{i},\pmb{i}+\pmb{x}}=\chi\tau^{1}+\eta\tau^{2}, & & u_{\pmb{i},\pmb{i}+\pmb{y}}=(-)^{i_{x}}(\chi\tau^{1}-\eta\tau^{2}),\\[2pt]
u_{\pmb{i},\pmb{i}+2\pmb{x}+\pmb{y}}=(-)^{i_{x}}\lambda\tau^{3}, & & u_{\pmb{i},\pmb{i}-\pmb{x}+2\pmb{y}}=-\lambda\tau^{3},\\[2pt]
u_{\pmb{i},\pmb{i}+2\pmb{x}-\pmb{y}}=(-)^{i_{x}}\lambda\tau^{3}, & & u_{\pmb{i},\pmb{i}+\pmb{x}+2\pmb{y}}=-\lambda\tau^{3},\\[2pt]
a_{0}^{l}=0, & &
\end{array}
\tag{9.6.7}
\end{equation*}

以及 Z2Bzz1n：
\begin{equation*}
\begin{array}{rcl}
u_{\pmb{i},\pmb{i}+\pmb{x}}=\chi\tau^{1}+\eta\tau^{2}, & & u_{\pmb{i},\pmb{i}+\pmb{y}}=(-)^{i_{x}}(\chi\tau^{1}-\eta\tau^{2}),\\[2pt]
u_{\pmb{i},\pmb{i}+2\pmb{x}+\pmb{y}}=(-)^{i_{x}}(\chi_{1}\tau^{1}+\eta_{1}\tau^{2}+\lambda\tau^{3}), & & u_{\pmb{i},\pmb{i}-\pmb{x}+2\pmb{y}}=\chi_{1}\tau^{1}-\eta_{1}\tau^{2}+\lambda\tau^{3},\\[2pt]
u_{\pmb{i},\pmb{i}+2\pmb{x}-\pmb{y}}=(-)^{i_{x}}(\chi_{1}\tau^{1}+\eta_{1}\tau^{2}-\lambda\tau^{3}), & & u_{\pmb{i},\pmb{i}+\pmb{x}+2\pmb{y}}=\chi_{1}\tau^{1}-\eta_{1}\tau^{2}-\lambda\tau^{3},\\[2pt]
a_{0}^{l}=0. & &
\end{array}
\tag{9.6.8}
\end{equation*}

对于式 (9.6.1)、(9.6.2)、(9.6.3) 与 (9.6.4) 给出的前四个 $Z_{2}$ 自旋液体（假定式 (9.6.1) 中 $a_{0}^{1}$ 较小），自旋子在四个孤立点处无能隙，并具有线性色散（见图 9.8）。因此，这四个拟设描述对称的 $Z_{2}$ 线性自旋液体。第二个 $Z_{2}$ 自旋液体 Z2Azz13 的单自旋子（single-spinon）色散相当有趣：它不具有 $90^{\circ}$ 旋转对称性。这与基态中的 $90^{\circ}$ 对称性并不矛盾，因为奇数个自旋子的激发永远无法满足约束，因而是被禁止的。自旋子色散具有绕 $\pmb{k}=(0,\pi)$ 的 $90^{\circ}$ 旋转对称性，而偶数个自旋子的激发谱则具有 $90^{\circ}$ 旋转对称性。

%FIG{9.8}: {$Z_{2}$ 线性自旋液体的自旋子色散 $E_{+}(\pmb{k})$ 作为 $(k_{x}/2\pi,\,k_{y}/2\pi)$ 的函数的等高线图：(a) 式 (9.6.1) 的 Z2A0013 态；(b) 式 (9.6.2) 的 Z2Azz13 态；(c) 式 (9.6.3) 的 Z2A001n 态；(d) 式 (9.6.4) 的 Z2Azz1n 态。}

我们想指出，当 $a_{0}^{1}$ 较大时，式 (9.6.1) 可能有有能隙的自旋子谱。因此，Z2A0013 态既可能是有能隙自旋液体，也可能是无能隙自旋液体。其余三个 Z2A 态总是无能隙的，是 $Z_{2}$ 线性态。在 9.9 节中我们将证明，所有 $Z_{2}$ 有能隙态与 $Z_{2}$ 线性态都是稳定的。这些态的平均场结果可以应用于物理自旋液体。

接下来让我们考虑式 (9.6.5) 中的拟设 Z2B0013。拟设 (9.6.5) 的自旋子谱由下式的本征值确定：
\begin{equation*}
H=-2\big[\chi\cos(k_{x})\Gamma_{0}-\eta\cos(k_{x})\Gamma_{2}-\chi\cos(k_{y})\Gamma_{1}+\eta\cos(k_{y})\Gamma_{3}\big]+\lambda\Gamma_{4}
\end{equation*}

%FIG{9.9}: {$Z_{2}$ 线性态的自旋子色散 $\min(E_{1}(\pmb{k}),E_{2}(\pmb{k}))$ 作为 $(k_{x}/2\pi,\,k_{y}/2\pi)$ 的函数的等高线图：(a) 式 (9.6.5) 的 Z2B0013 态；(b) 式 (9.6.6) 的 Z2Bzz13 态；(c) 式 (9.6.7) 的 Z2B001n 态；(d) 式 (9.6.8) 的 Z2Bzz1n 态。}

其中 $k_{x}\in(0,\pi)$，$k_{y}\in(-\pi,\pi)$，而
\begin{equation*}
\begin{array}{rcl}
\Gamma_{0}=\tau^{1}\otimes\tau^{3}, & \quad\Gamma_{1}=\tau^{1}\otimes\tau^{1}, & \quad\Gamma_{2}=\tau^{2}\otimes\tau^{3},\\[2pt]
\Gamma_{3}=\tau^{2}\otimes\tau^{1}, & \quad\Gamma_{4}=\tau^{1}\otimes\tau^{0}, &
\end{array}
\end{equation*}
这里已取 $\gamma_{1,2}=\lambda_{2}=0$。自旋子色散的四支能带具有 $\pm E_{1}(\pmb{k})$、$\pm E_{2}(\pmb{k})$ 的形式。我们发现自旋子谱在 $\pmb{k}=(\pi/2,\pm\pi/2)$ 附近的 8 个孤立点处为零（见图 9.9(a)）。因此，Z2B0013 态是一个 $Z_{2}$ 线性自旋液体。

其余三个 Z2B 态的自旋子谱可以用类似方法求出，谱图画在图 9.9 中。这三个 Z2B 态也是 $Z_{2}$ 线性态。无能隙自旋子只出现在 $\big(\frac{\pi}{2},\pm\frac{\pi}{2}\big)$ 处。

% ---------------- 新增术语（ch9_e） ----------------
% | staggered-flux phase | 交错通量相（ch9_c 已收 staggered-flux state 交错通量态） |
% | U(1)/SU(2) projective symmetry group | U(1)/SU(2) 投影对称群（= U(1)/SU(2) PSG） |
% | type-U1A/U1B/U1C PSG | U1A/U1B/U1C 型 PSG |
% | label (of a PSG) | 标记（PSG 的标记） |
% | quantum critical state | 量子临界态 |
% | critical line | 临界线 |
% | unstable fixed point | 不稳定不动点 |
% | destabilize | 使……失稳 |
% | single-spinon (excitation/spectrum) | 单自旋子（激发/谱） |
% | non-collinear flux | 非共线通量（collinear 共线的，术语表已收） |
% | contour plot | 等高线图 |
% | the zoo of ... | ……动物园（章节名修辞，照译） |
